\section{Introduction}
\label{sec:intro}

% ---- 脉络--------------------------------------------------------
% P1  World action model（WAM）：用预训练视频生成模型做 backbone，policy 继承世界知识与动态先验。
% P2  代价与矛盾：backbone 的 token 空间由生成目标决定（稠密、均匀的时空网格），
%     而 action head 只需要其中一小部分；推理成本跟着生成形状走，而不是跟着控制需要走。
% P3  为什么现成的 token 剪枝在 WAM 的部署点上不成立：
%     相似度 != 冗余（future 帧近似噪声占位）；删除会破坏 backbone 的结构；
%     空间与时间两个轴的位置约定相反。
% P4  观点：应由 action 而不是视觉内容决定 backbone 在哪里、以什么分辨率计算
%     （在常规的 video -> action 之外，加上 action -> video 的通路）。
% P5  方法一段：C1 action-aware router；C2 空间分辨率分配（合并而非删除）；
%     C3 router 引导的时间分辨率 + 紧凑时间位置；可 fold 的 decoder 适配。
% P6  结果一句（对 native 只能写"不可区分"，不能写"超过"）。
% 贡献列表：C1 / C2 / C3（分析部分作为动机，不另算）。
% 若用 teaser：main.tex 中有整栏 teaser 的写法，label 为 fig:teaser。
% ---------------------------------------------------------------------------

% ---- 数字来源与措辞边界（Claude 起草时记录）----------------------------------
% 框架：action-aware router，不写成 pruning；效率只作为副产品，延迟尚未测量（见文末 \TODO）。
% 只引用已写的 Sec. 4 与 fig:overview。Sec. 3 已删除，P3 与结果段的论据（相似度不优于随机、删除在小预算下崩溃）需在 Sec. 5 给出。
% 本节新增的 16 个 \cite 条目还没有加入 main.bib（按"只改 intro"的要求没有动 main.bib），加入之前这些引用显示为 [?]。
% LIBERO，每个套件 n=500，逐 episode 配对 exact McNemar；预算只作用于 16,800 个 future token，2,400 个 observation token 全部保留：
%   20% future（总 5,760 = 30%）：Spatial 441 vs 452（p=0.22），Object 470 vs 470（p=1），Goal 447 vs 451（p=0.71）-> 只写"不可区分"。
%     Object 的冻结 router 配方与 Spatial/Goal 不同（v2-distill），Sec. 5 需注明。
%   30% future（总 7,440）：92.0 / 96.4 / 90.2 vs 90.4 / 94.0 / 90.2；Object p=0.036，Holm 后 0.32 -> 同样只写"不可区分"。
%   10% future（总 4,080，Spatial）：354 = 70.8%；同一 router、各自重训 adapter 的 token 删除 93 = 18.6%、tube 删除 259 = 51.8%；
%     同一算子但段用原帧号作时间位置 275 = 55.0%（p=7.5e-13）。
%   E1（反转排序 + 按反转分配重训 adapter，Spatial）：20% 88.2 vs 66.2（-22.0pp），10% 70.8 vs 28.4（-42.4pp），Holm p<=6e-21。
%   E3（仿真分割，Spatial）：满分辨率组覆盖任务物体像素 90%、机械臂 84%、背景 26%，随机分配为 33%。
%     不能写"按指令选择"：同类干扰碗无差别，换指令后 96--97% 的满分辨率组不变。
%   E2（离线，Spatial）：router 的动作失真低于 similarity / motion / block-2 self-attention / edge / random（Holm p<=4e-4）。
%     不与 decoder cross-attention 比较，不写"优于 attention 类准则"。
%   "相似度不优于随机"：E2 中 similarity 的失真点估计高于 random；早期闭环（Object 6,000 tokens，n=100）随机排序 96 = 相似度排序 96。
%   adapter 可 fold：闭环核查是在旧算子上做的（合并 -1.10pp，p=0.15），定稿算子上尚未复核。
% 可选，是否纳入第二个 backbone 由你决定（DeVA 尚未达到 >= native，不建议进 intro）：
%   On a second WAM, DiT4DiT~\cite{ma2026dit4dit}, allocating the future tokens in the same way removes 24.9\% of all tokens on LIBERO
%   without changing the success rate over 1,500 paired episodes (98.9\% for both).
%   依据：589/784 tokens，合并 1483 vs 1483，9/9 不一致对。LIBERO 上只有 1 个 future slice，只验证了空间部分；
%   时间部分只在 RoboCasa-GR1 上验证（490/588 tokens，55.4% vs 54.9%，n=1,200），两部分没有在同一格中同时启用。
% ---------------------------------------------------------------------------

Video generation models trained on large-scale video learn how scenes look and how they evolve~\cite{nvidia2025cosmos,wan2025wan}.
World action models (WAMs) turn such models into robot policies~\cite{pai2025mimicvideo,hu2025vpp,kim2026cosmospolicy,ye2026dreamzero,ma2026dit4dit,zhang2026deva}.
In a common design~\cite{pai2025mimicvideo,ma2026dit4dit,zhang2026deva}, a video diffusion transformer (DiT)~\cite{peebles2023dit} predicts how the scene will evolve from the current observation and the instruction,
and an action decoder reads the backbone's intermediate representation of this predicted future to generate the robot's actions.
The policy thus inherits the visual and physical knowledge of the video model instead of learning it from robot demonstrations alone.

The representation the action decoder reads, however, keeps the layout of video generation.
Generating a video requires every region of every frame, so each future frame is a dense grid of tokens at uniform resolution;
in the mimic-video policy we build on~\cite{pai2025mimicvideo}, the 14 future frames alone account for 16,800 of the 19,200 tokens the backbone computes.
The action depends on much less: the robot, the objects it acts on and the space it moves through occupy only part of the scene.
Yet the backbone computes every region at full resolution, and the decoder attends to all of them, whether or not they matter for the action.
The representation, and with it the cost of inference, is shaped by what video generation needs rather than by what control needs.
This raises the question we study: which part of the predicted future does the action rely on, and where can the backbone compute at a coarser resolution without changing the action?

Token reduction methods address a related question with visual criteria.
They merge tokens that look alike~\cite{bolya2023tome,bolya2023tomesd,hyun2025sttm}, drop tokens that receive little attention~\cite{liang2022evit,chen2024fastv},
or, in vision-language-action models, reuse tokens that change little between control steps and keep those that attend to the instruction~\cite{xu2025vlacache,yang2025efficientvla}.
These criteria measure whether a token is redundant in appearance or salient to the model, not whether the action depends on it, and many of these methods discard the tokens they do not keep.
In a WAM, both choices fall short.
Choosing which regions keep full resolution by feature similarity is no better than choosing them at random,
and deleting tokens removes parts of the scene from the representation altogether, which degrades the policy sharply when few tokens are kept.

We argue that in a WAM, the action should decide where the backbone computes in detail.
We call a token action-relevant if reducing its resolution changes the predicted action.
The objective of the action decoder measures this change directly, so it can supervise where the resolution goes.
This adds a path from the action to the video backbone, alongside the usual path from the video to the action.

We propose \method, an action-aware router for WAMs (\cref{fig:overview}).
After an early block of the frozen video DiT, a lightweight router scores every future token by its relevance to the action (\cref{sec:router}).
The router is supervised only through the action decoder:
its labels come from the gradient of the decoder's objective with respect to the router's scores, and no term of its loss involves appearance, reconstruction or similarity between tokens.
The scores then decide at which resolution the remaining blocks compute each region, in space and in time.
Spatially, we group the tubes, the tokens at one grid position across the future frames, into $2{\times}2$ neighborhoods;
the most relevant groups keep full resolution, and every other group is merged into a single tube at its original grid position, so no region of the scene is removed (\cref{sec:spatial}).
Temporally, each tube is divided into a number of segments set by its relevance, and the segments take consecutive temporal positions, the pattern of consecutive frames the video model was trained on (\cref{sec:temporal}).
The video model and the action decoder stay frozen; only the router and a low-rank adapter~\cite{hu2022lora} in the decoder are trained, and the adapter can be folded into the decoder's weights after training.
Because the blocks after the router and the cross-attention of the decoder operate on the reduced sequence, the allocation also reduces computation.

On LIBERO~\cite{liu2023libero}, with 20\% of the future tokens kept (30\% of all tokens), \method is statistically indistinguishable from the full-context policy on all three suites
(Spatial 88.2\% vs.\ 90.4\%, Object 94.0\% vs.\ 94.0\%, Goal 89.4\% vs.\ 90.2\%; 500 paired episodes per suite).
With 10\% kept on LIBERO-Spatial, \method achieves 70.8\% success, whereas deleting tokens or whole tubes ranked by the same router, each with its own adapter, reaches 18.6\% and 51.8\%;
giving the segments their original frame positions instead of consecutive ones reaches 55.0\%.
Three analyses on LIBERO-Spatial indicate that the router captures what the action relies on.
Reversing its ranking at the same budget lowers success by 22.0 points at 20\% and by 42.4 points at 10\%, even when the adapter is retrained for the reversed allocation.
The full-resolution regions cover 90\% of the pixels of the task objects and 84\% of those of the robot, but only 26\% of the background, where a random allocation would cover 33\%.
Offline, the router's choice of full-resolution regions distorts the predicted action less than choosing them by feature similarity, motion, early-layer self-attention or edge energy, or at random.
\TODO{Efficiency: latency of the deployed budgets, once measured.}

In summary, our contributions are as follows.
\begin{itemize}
  \item We propose an action-aware router for WAMs that is supervised only through the gradient of the action decoder's objective,
  and show that the regions it keeps at full resolution are necessary for the policy and cover the robot and the task objects.
  \item We allocate spatial resolution by action relevance: the most relevant regions keep full resolution and the others are merged rather than deleted,
  with every token kept at its original grid position so that it stays aligned with the observation.
  \item We allocate temporal resolution by action relevance: each region receives a number of time segments set by its relevance, and the segments take consecutive temporal positions.
\end{itemize}
